\documentclass[11pt,a4paper]{article}
\usepackage[T1]{fontenc}
\usepackage[utf8]{inputenc}
\usepackage{amsmath,amssymb,amsthm}
\usepackage[margin=1in]{geometry}
\usepackage[colorlinks=true,linkcolor=blue,urlcolor=blue]{hyperref}
\theoremstyle{plain}
\newtheorem{theorem}{Theorem}
\newtheorem{lemma}[theorem]{Lemma}
\newtheorem{proposition}[theorem]{Proposition}
\newtheorem{corollary}[theorem]{Corollary}
\theoremstyle{definition}
\newtheorem{remark}[theorem]{Remark}

\title{The empirical recurrence for OEIS A237144, proved by transfer matrix}
\author{Adrian Perez Fontelles\\ \small Independent researcher}
\date{7 September 2026}
\begin{document}
\maketitle

\begin{abstract}
OEIS A237144 carries an empirical recurrence of order $70$ contributed by R. H. Hardin. It is
true, and it is decidable rather than empirical. The entry's condition is local to a bounded window of consecutive lines, so
the admissible configurations of such a window are the vertices of a finite digraph and the
arrays counted are exactly the walks in it; hence $a(n)$ is a walk count on $S=216$ vertices
and is $C$-finite. Whether the conjectured recurrence annihilates it is then settled by a
finite exact computation, with the Cayley--Hamilton theorem bounding the work.
\end{abstract}

\noindent\small 2020 Mathematics Subject Classification. 05A15, 05B45, 15B36.\normalsize

\section{The conjecture}

OEIS A237144 is ``Number of (n+1)X(2+1) 0..3 arrays with the maximum plus the lower median of every 2X2 subblock equal.''. It has offset $1$ and begins
\[
1472, 17616, 197724, 2463178, 30395540, 386318536, 4907807434, 63050866982,\ \dots
\]
The entry states, as a conjecture contributed by its author and never marked settled:
\begin{quote}
Empirical recurrence of order 70 (see link above)
\end{quote}
As of the ``Last modified'' line on the live entry (Jul 23 2025 09:53:12, revision 7) this is still
recorded as empirical, and nothing on the entry records it as proved.

\section{The count is a walk count}

The condition the entry imposes is local: it constrains a bounded window of consecutive lines and nothing beyond. Take as
vertices the admissible configurations of one such window and put an edge from $u$ to $v$
whenever $v$ may follow $u$, which the condition decides by inspection of the two. An
admissible array is then exactly a walk, so with $M$ the adjacency matrix on $S=216$ vertices
and $\iota,\tau$ the vectors recording which windows may begin and end an array,
\[
a(n)\;=\;\iota^{\!\top}M^{\,n}\tau ,
\]
the exponent fixed by the entry's own shape and checked against its published terms. In
particular $a$ is $C$-finite.

\section{The proof}

\begin{lemma}\label{lem:crit}
Let $q(t)=t^{70}-\sum_i c_it^{70-i}$ be the characteristic polynomial of the
conjectured recurrence. The residual $a(n)-\sum_ic_ia(n-i)$ equals
$\iota^{\!\top}M^{\,j}q(M)\tau$ for the corresponding walk index $j$, so the recurrence
holds from some point on if and only if $\iota^{\!\top}M^{j}q(M)\tau=0$ for all large $j$,
and it is enough to examine $j<S$.
\end{lemma}

\begin{proof}
Factor $M^{j}$ out of $M^{j+70}-\sum_ic_iM^{j+70-i}$. For the bound, $M^{S}$
is an integer combination of $I,M,\dots,M^{S-1}$ by Cayley--Hamilton, so the numbers
$u_j=\iota^{\!\top}M^{j}q(M)\tau$ obey the monic recurrence given by the characteristic
polynomial of $M$; $S$ consecutive zeros force every later one, and the last nonzero $u_j$
pins the threshold exactly.
\end{proof}

\begin{theorem}
\[
a(n)\;=\;39\,a(n-1) - 160\,a(n-2) - 12044\,a(n-3) + 145902\,a(n-4) + 1301357\,a(n-5) - 28998486\,a(n-6) - 23143247\,a(n-7) + 2983103077\,a(n-8) - 8372689257\,a(n-9) - 186522916060\,a(n-10) + 1021940846958\,a(n-11) + 7344299435320\,a(n-12) - 63749321988945\,a(n-13) - 166330077893225\,a(n-14) + 2599781261562279\,a(n-15) + 703737238969351\,a(n-16) - 74689355357780259\,a(n-17) + 96224930618884890\,a(n-18) + 1554555314427677603\,a(n-19) - 4037104924465079240\,a(n-20) - 23439383477887989510\,a(n-21) + 94468254131551344798\,a(n-22) + 245986355248174699197\,a(n-23) - 1538068456145135501031\,a(n-24) - 1509409575078360285981\,a(n-25) + 18585206013649610854908\,a(n-26) - 969679345185840587997\,a(n-27) - 171019089996424356270276\,a(n-28) + 141770311292210122719180\,a(n-29) + 1208287224153788793868230\,a(n-30) - 1833180013932432304366755\,a(n-31) - 6515970743144530671621596\,a(n-32) + 14722519233845227915335344\,a(n-33) + 26100016339490264593135930\,a(n-34) - 85911886883585776474149659\,a(n-35) - 71381602952343952813346574\,a(n-36) + 382888187578277599439919469\,a(n-37) + 88741065470822620977868868\,a(n-38) - 1330364967974923022020690482\,a(n-39) + 263612103359188023210204998\,a(n-40) + 3632822012287157342790093062\,a(n-41) - 1959772874708864128893661132\,a(n-42) - 7800203620907892204716451725\,a(n-43) + 6553800746731419550738384340\,a(n-44) + 13090078467385195437896000593\,a(n-45) - 14918802129637674735916961289\,a(n-46) - 16929578769554397337783200256\,a(n-47) + 25163812912874646295081404754\,a(n-48) + 16411352874801018496786413612\,a(n-49) - 32344378109509927413894038376\,a(n-50) - 11228174964373034514526126784\,a(n-51) + 31947787660331753005976517728\,a(n-52) + 4525159837076488065598251776\,a(n-53) - 24213200715632083926698812032\,a(n-54) + 13921824249548065145684992\,a(n-55) + 13966608945681628925157484544\,a(n-56) - 1371967094177676433883704320\,a(n-57) - 6046924385419647588041842688\,a(n-58) + 976284382984348699783348224\,a(n-59) + 1926631847549153169292263424\,a(n-60) - 376832937576072420213325824\,a(n-61) - 439760960432071967326076928\,a(n-62) + 89346900076942136064344064\,a(n-63) + 69296632458273967331868672\,a(n-64) - 12830845605765831928578048\,a(n-65) - 7130511556499570128060416\,a(n-66) + 1015785531774837102477312\,a(n-67) + 432510443602795258970112\,a(n-68) - 33789326924988071018496\,a(n-69) - 11836509492045108215808\,a(n-70)
\]
for every $n>70$.
\end{theorem}

\begin{proof}
The vector $w=q(M)\tau$ was formed by $70$ matrix--vector products in exact integer
arithmetic and $\iota^{\!\top}M^{j}w$ evaluated for $j=0,1,\dots$, again exactly; those
integers vanish from the index corresponding to $n=70+1$ onwards and the run continued
until the working vector was identically zero, which settles every larger $j$ at once.
\end{proof}

\section{Verification}

Three checks, each independent of the algebra above.

First, the model reproduces all $15$ terms the entry publishes, exactly, in integer
arithmetic. This is what ties the model to the entry: the digraph is built by reading the
entry's English, and a misreading gives different counts.

Second, the conjectured recurrence was evaluated directly on those published terms, with no
matrices involved, and holds wherever the proved range and the data overlap.

Third, the annihilation test of Lemma~\ref{lem:crit} was carried out in exact integer
arithmetic on the whole vector rather than on a sampled prefix.

\begin{thebibliography}{9}
\bibitem{oeis} The OEIS Foundation, \emph{The On-Line Encyclopedia of Integer
Sequences}, \url{https://oeis.org}, 2026. Sequence A237144.
\bibitem{stanley} R.~P.~Stanley, \emph{Enumerative Combinatorics}, Volume 1, 2nd ed.,
Cambridge University Press, 2011. (Transfer-matrix method, Section 4.7.)
\bibitem{fs} P.~Flajolet and R.~Sedgewick, \emph{Analytic Combinatorics}, Cambridge
University Press, 2009.
\bibitem{horn} R.~A.~Horn and C.~R.~Johnson, \emph{Matrix Analysis}, 2nd ed., Cambridge
University Press, 2013.
\end{thebibliography}

\end{document}
